\scanpage{12}
i.e la donnée d’une situation $X'_i\mathbin{-}X_i$ sur $k_i$, \deleted{dont} et un iso. $(\unsure{\mathscr X|\mathcal X|\mathfrak X}'\mathbin{-}\mathscr X)\simeq(X'_i\mathbin{-}X_i)\otimes_{k_i}\bar k$, de ceci \unsure{se|on} déduit une extension de schémas en groupes finis étales sur $k_i$
\diagram{(27)\quad\begin{tikzcd}[ampersand replacement=\&,column sep=small]
1\arrow[r,dash]\&G_i\arrow[r,dash]\&\unsure{\mathfrak G|G|\mathscr G}'_i\arrow[r]\&\mathfrak G_i\arrow[r,dash]\&1
\end{tikzcd}}
et on exige que $\mathfrak G'_i$ soit un groupe constant [i.e. que les autom ``\unsure{géom|globaux|galois|géom.}'' de $\mathscr X'\mathbin{-}\mathscr X$ (sur $\bar k$) \deleted{\unsure{soient}} ``définis sur $k_i$'' \dots]

Il faut bien noter que pour $k_i$ fixé, il peut y avoir plusieurs données non isomorphes de cette espèce sur $k_i$, i.e. plusieurs scindages partiels $\Gamma_i\mathbin{-}\unsure{\mathcal E|\mathscr E}'$ qui centralisent $\mathfrak G'$. \deleted{\unsure{se|La}} La propriété d’unicité \deleted{consiste} évidente consiste \unsure{que|pour} tt \unsure{s’égalisent|s’égaliser}, quitte à passer à une sous-extension finie assez grande ....

De façon précise, introduisons
\[
\unsure{\mathcal Z|Z|\mathscr Z}=\mathrm{Cent}_{\unsure{\mathcal G|\mathscr G}'}(\mathfrak G')\tag*{(28)}
\]
% archive sheet number: 11
